\documentclass[12pt]{article}
\usepackage{amsmath, amssymb, physics}
\usepackage{geometry}
\usepackage{hyperref}
\geometry{a4paper, margin=1in}

\title{Physics and Mathematics of Holland-Burnett State Initialization and Metrology Metrics}
\author{A Theoretical Review}
\date{\today}

\begin{document}
\maketitle

\section*{Overview}
This notebook structure evaluates Holland-Burnett (HB) states tracking Quantum Fisher Information (QFI) against absolute physical boundaries over specific noisy parameters. Our evaluation maps limits of $N$ photon signals subjected to $d$ phases, computing performance drops over probabilistic shared vs. independent noise distributions ($p, \eta$) effectively.

---

\section*{Cell 1: Imports and Dependencies}
Initializes matrix arrays efficiently utilizing Pandas mapping and \texttt{qutip} frameworks dynamically. Necessary tools including \texttt{tensor}, \texttt{qeye}, \texttt{basis}, and \texttt{ket2dm} set up vector mathematics tracking quantum evolutions natively instead of analytically guessing structures.

\section*{Cell 2: The Core Physics Engine}

\subsection*{2.1 Quantum Fisher Information Matrix (QFIM)}
The Symmetric Logarithmic Derivative (SLD) $L_{\mu}$ governs the matrix geometry necessary to derive Cramer-Rao boundaries (CRB). For a state $\rho$ subject to phase displacement estimations $U(\boldsymbol{\theta}) = \exp(-i \sum_{\mu} \hat{n}_{\mu} \theta_{\mu})$, the partial derivative evaluated at the phase origin implies:
\begin{equation}
    \partial_\mu \rho \Big|_{\boldsymbol{\theta} = 0} = -i [\hat{n}_\mu, \rho]
\end{equation}
To calculate this practically, we resolve $\rho$ into its spectral eigenbasis: $\rho = \sum_k \lambda_k |k\rangle\langle k|$. Consequently, the partial derivative map transitions analytically to:
\begin{equation}
    \langle k | \partial_\mu \rho | m \rangle = -i(\lambda_m - \lambda_k)\langle k | \hat{n}_\mu | m \rangle
\end{equation}
In our codebase (seen precisely in \texttt{compute\_rho\_derivative}), this manifests by establishing matrices for eigenvalue variations evaluated natively. Proceeding forward, the QFIM bridges from the derivatives:
\begin{equation}
    F_{\mu\nu} = \operatorname{Re}\left(  \sum_{\lambda_m + \lambda_k > 0} \frac{2}{\lambda_n + \lambda_m} \langle k| \partial_\mu \rho |m\rangle \langle m| \partial_\nu \rho | k\rangle \right)
\end{equation}
This bounds identically the multiple independent phases' ultimate precision evaluated as standard Cramer-Rao variance:
\begin{equation}
    \Delta^2 (\hat{\boldsymbol{\theta}}) \geq \operatorname{Tr}\left(\mathcal{F}^{-1}_{\boldsymbol{\theta}}\right)
\end{equation}

\subsection*{2.2 The Noise Channel ($\mathcal{N}_{\text{tot}}$)}
The beam-splitter simulation models practical photonic damping (loss mechanism mapping to non-unity transmissivity $\eta$). The Stinespring dilation interaction linking our active signal modes $\hat{a}_S$ against surrounding reservoir vacuums $\hat{a}_E$:
\begin{equation}
    \hat{V} = \exp\left[\theta_{\text{int}}\left( \hat{a}_S^\dagger \hat{a}_E - \hat{a}_S \hat{a}_E^\dagger \right)\right] 
\end{equation}
Under standard transformations, mapping relates mathematically to evaluating $\cos^2(\theta_{\text{int}}) = \eta$, ensuring $\theta_{\text{int}} = \arccos(\sqrt{\eta})$. 

The script parses this globally against two distinct decoherence profiles linked directly to a mixing probability $p$:
\begin{itemize}
    \item \textbf{Independent Error ($\mathcal{N}_{\text{free}}$)}: Distributes exclusive independent vacuum ancillary channels against sequentially unlinked operators per each signal register $i \in K$.
    \item \textbf{Correlated Noise ($\mathcal{N}_{\text{int}}$)}: Operates globally over synchronized phase loss wherein all parallel phase registers route unitarily against a single shared environment entity.
\end{itemize}

\section*{Cell 3: State Initialization Equations}

\subsection*{3.1 The Holland-Burnett State ($|\psi_{HB}\rangle$)}
Standard initialization bounds map $N$ resources globally over $d$ phases, extending fundamentally into $d+1$ total components (registering signal versus reference spaces). 
The $HB(n, d)$ architecture explicitly deposits identical particle distributions $n=N/(d+1)$ targeting respective ports of a multiport Quantum Fourier Transform (QFT) interferometer gate.
We construct dynamic transformed creation structures via:
\begin{equation}
    \hat{b}_k^{\dagger} = \frac{1}{\sqrt{d+1}}\sum_{j=0}^d e^{i\frac{2\pi}{d+1}j \cdot k} \, \hat{a}_j^{\dagger}
\end{equation}
Resulting mathematically over vacuum applications:
\begin{equation}
    |\psi_{HB}\rangle \propto (\hat{b}_0^\dagger)^n (\hat{b}_1^\dagger)^n \dots (\hat{b}_d^\dagger)^n |0\rangle^{\otimes d+1} 
\end{equation}
Resolving specifically $N=3$, $d=2$ yields mathematically the output proxy matching purely derived state symmetries expanding algebraically to distributions resolving exclusively the limits.

\subsection*{3.2 The Reference Bound State ($|\psi_s\rangle$)}
Evaluating directly the theoretical peak established natively in Humphreys Equation 8 optimally minimizing variance bounding Equation 10 limit:
\begin{equation}
    |\psi_S\rangle = \beta|0_{\text{sig}}\rangle \otimes |N_{\text{ref}}\rangle + \alpha \sum_{j=1}^d |N^{(j)}_{\text{sig}}\rangle \otimes |0_{\text{ref}}\rangle
\end{equation}
Where mathematically evaluating the scale amplitude yields exact derivations isolating optimal components:
\begin{equation}
    \alpha = \frac{1}{\sqrt{d+\sqrt{d}}} \quad \text{and} \quad \beta = \frac{1}{\sqrt{1+\sqrt{d}}}
\end{equation}
These arrays pass unitarily directly as exact numerical kets resolving fidelity approximations evaluated inside our matrices computationally.

\section*{Cell 4: Metric Evaluation and Layout Logic}
This section handles plotting matrices formatting Cramer-Rao boundaries natively against $p$ scales. 
Instead of optimizing, we plot identically our statically parsed states over grids. 
We establish analytical values:
\begin{itemize}
    \item \textbf{Standard Quantum Limit (SQL)} scales $\approx d^2 / N$ for total combined variance scales efficiently.
    \item \textbf{Heisenberg NOON Bound} scales $\approx d^3 / N^2$. 
    \item \textbf{Humphreys Ideal Multipixel Eq 10 Limit} natively maps $\approx \frac{d(1+\sqrt{d})^2}{4N^2}$.
\end{itemize}

\section*{Cell 5: Global Execution Block}
Combines states from Cell 3, evaluates grids from Cell 2 traces spanning ranges of independent vs shared mixtures ($p$) and tracks plotting maps from Cell 4 automatically to visualize exactly evaluated footprints. Initial calculations compute exactly $Fidity = |\langle \psi_{HB} | \psi_{s} \rangle|^2 $.

\section*{Cell 6: QCRB Under Photon Loss — Humphreys Fig 3 Replication}

This section replicates the core results of Humphreys et al.\ Figure 3, evaluating the performance of HB states under realistic photon-loss noise modelled by the Kraus operator formalism.

\subsection*{6.1 Loss Model: Kraus Operators}
Each signal mode $j \in \{1, \ldots, d\}$ is subjected to an independent amplitude-damping channel with transmission probability $\eta$. The Kraus operators are:
\begin{equation}
    E_k = \sum_{n=k}^{D-1} \sqrt{\binom{n}{k}} \, \eta^{(n-k)/2} \, (1-\eta)^{k/2} \, |n-k\rangle\langle n|, \quad k = 0, 1, \ldots, D-1
\end{equation}
The output state after loss on all $d$ signal modes (mode 0, the reference, is lossless):
\begin{equation}
    \rho_{\text{out}} = \prod_{j=1}^{d} \mathcal{E}_j(\rho_{\text{pure}}), \quad \mathcal{E}_j(\rho) = \sum_{k=0}^{D-1} E_k^{(j)} \rho \, E_k^{(j)\dagger}
\end{equation}

\textbf{Note on Computational Efficiency:} While previous iterations (Project A) utilized a direct Stinespring dilation (explicitly modeling the environment as additional modes), this implementation transitions to the Kraus operator formalism. This approach is mathematically equivalent but avoids the exponential memory growth associated with adding physical ancilla modes. This optimization is critical for replicating Humphreys Fig 3, as it allows us to handle Hilbert space dimensions for $N=12$ and $N=18$ that would otherwise crash the system.

\subsection*{6.2 QCRB vs Total Photon Number $N$ (Fig 3A, $d=2$ fixed)}
We fix $d=2$ (estimating 2 phases simultaneously) and sweep $n = 1, 2, \ldots, 6$, giving $N = 3n = 3, 6, 9, 12, 15, 18$. The Hilbert space dimension is $(N+1)^3$, reaching a maximum of $19^3 = 6{,}859$ for $N=18$.

Loss parameters tested: $\eta \in \{1.0, 0.95, 0.9, 0.85\}$.

\textbf{Key observations:}
\begin{itemize}
    \item All curves decrease monotonically with $N$ (more photons $\rightarrow$ better precision).
    \item The gap between the lossless ($\eta = 1$) and noisy curves \emph{increases} with $N$, indicating noise has a compounding effect at higher photon numbers.
    \item HB states outperform N00N states for $N \gtrsim 6$ even in the lossless case.
    \item For large $N$, the noisy HB curves approach but do not cross the ideal optimal bound $\frac{d(1+\sqrt{d})^2}{4N^2}$.
\end{itemize}

\textbf{Note:} The original Humphreys figure uses $d=3$ (4 modes). We use $d=2$ (3 modes) because the $d=3$ case at $N=12$ requires eigendecomposing a $28{,}561 \times 28{,}561$ matrix, which exceeds feasible computation time on available hardware.

Output file: \texttt{QCRB\_vs\_N\_d2\_noise.png} (full range), \texttt{QCRB\_vs\_N\_d2\_zoomin.png} (large-$N$ detail).

\subsection*{6.3 QCRB vs Number of Phases $d$ (Fig 3B, $n=1$ fixed)}
We fix $n=1$ and sweep $d = 1, 2, 3, 4$, giving $N = d+1 = 2, 3, 4, 5$. The Hilbert space dimension is $(d+2)^{d+1}$, reaching $6^5 = 7{,}776$ for $d=4$.

Loss parameters tested: $\eta \in \{1.0, 0.95, 0.9, 0.85\}$.

\textbf{Key observations:}
\begin{itemize}
    \item The total variance increases approximately linearly with $d$, consistent with the additive nature of independent phase uncertainties.
    \item The spacing between $\eta$ curves grows with $d$, showing that noise impact compounds across more phases.
    \item HB states consistently outperform N00N states at all values of $d$.
    \item The pure-state HB curve lies close to but above the optimal ideal bound, confirming that HB states are near-optimal but not perfectly so.
\end{itemize}

Output file: \texttt{QCRB\_vs\_d\_noise.png}.

\end{document}
